% ===========================================================================
\chapter{Changing things safely}\label{ch:editing}

Now the agent gets to change files. How it does that is one of the most
important choices in the whole design, so let's look at the options first.

\section{Three ways to edit a file}

\textbf{Write the whole file.} The model sends the complete new text.
Simple, and for a new file it is the only way. For a change to an existing
file it is wasteful (a one-line fix to a 2,000-line file means writing
2,000 lines), and it is risky: in all that copying, the model may change
something else by accident, or leave something out, and nothing tells
you.

\textbf{Replace lines by number.} ``Replace lines 32 and 33 with \ldots''.
Short, but fragile. Line numbers go stale: after one edit adds a line,
every number below it is off by one, and the model must keep track. Get it
wrong and the edit lands in the wrong place, again silently.

\textbf{Replace text by text.} ``Replace \emph{this} text with \emph{that}
text.'' The model quotes what it wants to change, exactly, and what to put
in its place. This is what we will use, with one rule that makes it safe:
\emph{the old text must occur exactly once in the file}. If it doesn't
occur at all, the model has misremembered the file, and should read it
again. If it occurs more than once, the edit is ambiguous, and the model
should quote more of the lines around it. Either way the edit is refused,
the file is untouched, and the model is told why. An edit either does
exactly what was meant, or nothing.

That rule is \code{count_occurrences(text, old_text) = 1}, using the helper
from Chapter~\ref{ch:setup}. Most coding agents in use today edit this way,
for these reasons.

\section{The edit tool}

The description repeats the rule, because the model has to follow it to
succeed: quote the text exactly, without the line numbers that
\code{read_file} added. This tool changes things, so it says so, and will
need permission.

\begin{append}{tools.nex}
class Edit_File inherit Tool
create
  make(ws: Workspace) do
    super.make(ws)
  end

feature
  name(): String do
    result := "edit_file"
  end

  description(): String do
    result := "Change a file by replacing one piece of its text with another. "
      + "`old_text` must match exactly one place in the file, character for character, "
      + "including spaces and line breaks (without the line numbers read_file shows). "
      + "If it matches more than one place, include more of the surrounding lines. "
      + "Use this for every change to an existing file."
  end

  parameters(): Map[String, Any] do
    result := {"path": parameter("string", "The file, relative to the workspace."),
      "old_text": parameter("string", "The exact text to replace."),
      "new_text": parameter("string", "The text to put in its place.")}
  end

  required(): Array[String] do
    result := ["path", "old_text", "new_text"]
  end

  changes_things(): Boolean do
    result := true
  end

  describe(input: Json_Value): String do
    result := "edit_file " + string_argument(input, "path", "?")
  end
\end{append}

\code{run} checks the path, then the edit. If both are fine, it replaces
the one occurrence and writes the file. What it sends back is not just
``done'': it shows the changed lines, numbered, with two lines either
side. The model sees the result of its edit straight away, as you would in
an editor, and can tell if something is off (the indentation, say)
without another \code{read_file}.

\begin{continued}{tools.nex}
  run(input: Json_Value): Tool_Outcome do
    let path: String := string_argument(input, "path", "")
    let old_text: String := string_argument(input, "old_text", "")
    let new_text: String := string_argument(input, "new_text", "")
    let problem: String := path_problem(path, "file")
    let before: String := ""
    if problem = "" then
      before := workspace.read(path)
      problem := edit_problem(before, path, old_text, new_text)
    end
    if problem /= "" then
      result := create Tool_Outcome.error(problem)
    else
      let after: String := replace_first(before, old_text, new_text)
      workspace.write(path, after)
      result := create Tool_Outcome.ok("Edited " + path + ": replaced "
        + count_text(lines_of(old_text).length, "line") + " with "
        + count_text(lines_of(new_text).length, "line") + ". It now reads:\n"
        + around(after, before.index_of(old_text), new_text))
    end
  end
\end{continued}

The checks get their own routine. Each refusal says what to do instead.
An empty \code{old_text} would match everywhere, and an edit that changes
nothing is a sign the model is confused; both are refused too.

\begin{continued}{tools.nex}
  -- Why `old_text` can't be replaced in `text`, or "" if it can.
  edit_problem(text, path, old_text, new_text: String): String do
    result := ""
    if old_text = "" then
      result := "`old_text` is empty; to write a whole file, use write_file"
    elseif old_text = new_text then
      result := "`old_text` and `new_text` are the same, so nothing would change"
    else
      let n: Integer := count_occurrences(text, old_text)
      if n = 0 then
        result := "`old_text` does not occur in " + path + ". Read the file again "
          + "and copy the text exactly, including spaces and line breaks."
      elseif n > 1 then
        result := "`old_text` occurs " + n.to_string + " times in " + path
          + ". Include more of the lines around it, so that it matches exactly one place."
      end
    end
  end
\end{continued}

To show the lines around the edit, \code{around} works out which line the
new text starts on (one more than the number of line breaks before it)
and how many lines it has:

\begin{continued}{tools.nex}
  private feature
    -- The lines of `text` around `new_text`, which starts at `start`,
    -- numbered, with two lines of context either side.
    around(text: String, start: Integer, new_text: String): String do
      let lines: Array[String] := lines_of(text)
      let first: Integer := count_occurrences(text.substring(0, start), "\n") + 1
      let last: Integer := first + lines_of(new_text).length - 1
      let from_line: Integer := (first - 2).max(1)
      let to_line: Integer := (last + 2).min(lines.length)
      result := numbered(lines.slice(from_line - 1, to_line), from_line)
    end
end
\end{continued}

\section{Writing whole files}

For new files, there is \code{write_file}. Its description steers the
model away from using it on existing files, for the reasons above.

\begin{append}{tools.nex}
class Write_File inherit Tool
create
  make(ws: Workspace) do
    super.make(ws)
  end

feature
  name(): String do
    result := "write_file"
  end

  description(): String do
    result := "Create a new file, or replace a whole file, with the given text. "
      + "Folders on the way are created. Use edit_file instead to change part of an existing file."
  end

  parameters(): Map[String, Any] do
    result := {"path": parameter("string", "The file, relative to the workspace."),
      "content": parameter("string", "The complete text of the file.")}
  end

  required(): Array[String] do
    result := ["path", "content"]
  end

  changes_things(): Boolean do
    result := true
  end

  describe(input: Json_Value): String do
    result := "write_file " + string_argument(input, "path", "?")
  end

  run(input: Json_Value): Tool_Outcome do
    let path: String := string_argument(input, "path", "")
    let content: String := string_argument(input, "content", "")
    let problem: String := path_problem(path, "any")
    if problem = "" and workspace.is_folder(path) then
      problem := "`" + path + "` is a folder"
    end
    if problem /= "" then
      result := create Tool_Outcome.error(problem)
    else
      let existed: Boolean := workspace.exists(path)
      workspace.write(path, content)
      result := create Tool_Outcome.ok((when existed then "Replaced " else "Created " end)
        + path + " (" + count_text(lines_of(content).length, "line") + ")")
    end
  end
end
\end{append}

The result says whether the file is new or replaced: if the model meant to
create a file and replaced one, that is something it should notice.

Let's try the edits a model might make on the sample project: one that is
ambiguous, one that quotes the file wrongly (it has the line numbers in),
and the real fix:

\begin{tryit}{try_edit.nex}
intern Tools

let ws: Workspace := fresh_workspace(user_path("sample_project"), user_path("scratch"))
let edit: Edit_File := create Edit_File.make(ws)

function show(outcome: Tool_Outcome) do
  console.print_line((when outcome.is_error then "error: " else "" end) + outcome.text)
end

show(edit.run(json_value('{"path": "loans.nex", "old_text": "result := ", "new_text": "result = "}')))
show(edit.run(json_value('{"path": "loans.nex", "old_text": "   32      if day >= due_on() then", "new_text": "if day > due_on() then"}')))
show(edit.run(json_value('{"path": "loans.nex", "old_text": "    if day >= due_on() then\n      result := day - due_on() + 1", "new_text": "    if day > due_on() then\n      result := day - due_on()"}')))
show(create Write_File.make(ws).run(json_value('{"path": "notes/fix.md", "content": "Late only after the due day.\n"}')))
\end{tryit}
\tryoutput{try_edit}

The two bad edits changed nothing. The good one shows its result. The
sample project is now fixed, in this scratch copy at least; we can't run
its checks yet, because the agent can't run commands. That is next.
